\section*{Chapter \ref{ch:b2:elements-in-the-earth} --- Elements in the Earth}

\begin{solution}{exo:b2:elements-in-the-earth:1}
Potassium is $2\times39.10/(2\times39.10+16.00) = 78.20/94.20 = 0.8301$ of the mass of
\ce{K2O}, so $\num{5130}\times0.8301 = \qty{4259}{ppm}$ of K. Titanium is
$47.87/79.87 = 0.5993$ of \ce{TiO2}: $\num{27310}\times0.5993 = \qty{16370}{ppm}$ of Ti.
\end{solution}

\begin{solution}{exo:b2:elements-in-the-earth:2}
Aluminium is lithophile and refractory: the silicate Earth kept all of it, in near
chondritic proportions (1.16). Sodium is lithophile but volatile (\qty{958}{K}): it was
lost before or during accretion, leaving 0.22. Nickel condenses with magnesium
(\qty{1353}{K}), so volatility cannot explain 0.077: it is siderophile, and most of it
went into the core. Iridium is refractory (\qty{1603}{K}) and strongly siderophile:
0.0029, nearly all of it in the core.
\end{solution}

\begin{solution}{exo:b2:elements-in-the-earth:3}
$\mathrm{La}_N = 31/0.232 = 133.6$, $\mathrm{Yb}_N = 2.0/0.163 = 12.27$, and
$(\mathrm{La}/\mathrm{Yb})_N = 10.9$. The crust is enriched in all the rare earths, and
eleven times more in the light ones than in the heavy ones: it is made of melts, which
take the incompatible light rare earths most readily.
\end{solution}

\begin{solution}{exo:b2:elements-in-the-earth:4}
$\mathrm{Sm}_N = 4.7/0.145 = 32.4$, $\mathrm{Eu}_N = 1.0/0.0546 = 18.3$,
$\mathrm{Gd}_N = 4.0/0.198 = 20.2$; $\mathrm{Eu}^* = \sqrt{32.4\times20.2} = 25.6$ and
$\mathrm{Eu}/\mathrm{Eu}^* = 18.3/25.6 = 0.72$, a negative anomaly. Plagioclase, which
takes in \ce{Eu^2+}, crystallised from the melts that built the upper crust and stayed
deeper down.
\end{solution}

\begin{solution}{exo:b2:elements-in-the-earth:5}
$\bar D_{\mathrm{La}} = 0.65\times0.058+0.18\times0.001+0.17\times0.001 = 0.0377 +
0.0002 + 0.0002 = 0.038$; $\bar D_{\mathrm{Lu}} = 0.3445+1.044+0.0306 = 1.42$. Lutetium is
compatible ($\bar D > 1$), lanthanum strongly incompatible. Garnet carries
$1.044/1.42 = 74$ per cent of the rock's lutetium, though it is only 18 per cent of the
rock.
\end{solution}

\begin{solution}{exo:b2:elements-in-the-earth:6}
For sodium, the silicate Earth's Na/Mg is $2670/\num{228000} = 0.01171$ and the CI
value $5010/\num{95870} = 0.05226$: the ratio is 0.224. For nickel, $1960/\num{228000}
= 0.008596$ against $\num{10640}/\num{95870} = 0.1110$: 0.0775. Sodium is
depleted by its volatility (\qty{958}{K}); nickel condenses above \qty{1300}{K} and is
three times more depleted than sodium because it is siderophile: the core took most of
it.
\end{solution}

\begin{solution}{exo:b2:elements-in-the-earth:7}
La: $5.21\times0.058 = \qty{0.302}{ppm}$; Lu: $0.53\times0.53 = \qty{0.281}{ppm}$. The
normalising values cancel in the ratio of ratios:
$(\mathrm{La}/\mathrm{Lu})_{N,\mathrm{cpx}}/(\mathrm{La}/\mathrm{Lu})_{N,\mathrm{melt}} =
D_{\mathrm{La}}/D_{\mathrm{Lu}} = 0.058/0.53 = 0.11$: the crystal's pattern is nine
times steeper towards the heavy rare earths than the melt's.
\end{solution}

\begin{solution}{exo:b2:elements-in-the-earth:8}
Crust: $(12/0.658)/(31/0.648) = 18.24/47.84 = 0.38$. Hawaiian basalt:
$(18.10/0.658)/(15.20/0.648) = 27.51/23.46 = 1.17$. MORB: $7.96/8.04 = 0.99$. Only the
crust, with a strong deficit of niobium, carries the subduction-zone signature.
\end{solution}

\begin{solution}{exo:b2:elements-in-the-earth:9}
\omterm{def:b2:elements-in-the-earth:bse}{Primitive mantle}: $240/0.0203 = \num{11800}$. MORB: K $= 0.160\times0.830\times10^4 =
\qty{1328}{ppm}$, K/U $= \num{11200}$. Crust: K $= 2.80\times0.830\times10^4 =
\qty{23240}{ppm}$, K/U $= \num{8600}$. Potassium rises a hundredfold from mantle to crust,
yet K/U changes by less than a third: the two elements are about equally incompatible,
and melting hardly separates them. Uranium is refractory, so its abundance follows from
the chondritic model; K/U measured on lavas then gives the planet's potassium, which,
being volatile, cannot be taken from \omterm{def:b2:planet-formation-meteorites:chondrule}{chondrites}.
\end{solution}

\begin{solution}{exo:b2:elements-in-the-earth:10}
At equilibrium $C_{\mathrm{rock}} = \bar D\,C_{\mathrm{melt}}$ for each element (and a
small melt fraction leaves the rock's composition nearly unchanged), so
$(C_1/C_2)_{\mathrm{melt}} = (C_{1}/\bar D_1)/(C_{2}/\bar D_2)$ for the rock's $C_1,
C_2$, and the ratio is $\bar D_2/\bar D_1$. For La and Lu: $1.419/0.03805 = 37$. The melt
of a garnet-bearing mantle is 37 times richer in La relative to Lu than its source: garnet
held back the heavy rare earths, and the melt's pattern is steep.
\end{solution}

\begin{solution}{exo:b2:elements-in-the-earth:11}
$C_{\mathrm{mix}} = fC_U + (1-f)C_M$ for each element: linear in $f$. The ratio
$\bigl(f\,\mathrm{Th}_U + (1-f)\mathrm{Th}_M\bigr)/\bigl(f\,\mathrm{Nb}_U +
(1-f)\mathrm{Nb}_M\bigr)$ is a quotient of two linear functions, a hyperbola in $f$.
Th/Nb $= 0.2$ when $f(10.5-2.4) + (1-f)(0.404-1.048) = 0$, that is $8.1f = 0.644(1-f)$,
$f = 0.644/8.744 = 0.074$: about 7 per cent of crust raises Th/Nb from 0.077 to 0.2. For
$f = 0.1$: Th $= 1.05 + 0.364 = 1.414$, Nb $= 1.2 + 4.716 = 5.916$, Th/Nb $= 0.24$.
\end{solution}

\begin{solution}{exo:b2:elements-in-the-earth:12}
Removing a component that holds neither element leaves both masses, hence their ratio,
unchanged. If all the aluminium was kept, $(\mathrm{Al}/\mathrm{Mg}) = 1.16$ times CI
means that the silicate Earth holds $1/1.16 = 0.86$ of its chondritic magnesium: 14 per
cent is missing. Silicon relative to aluminium is $0.83/1.16 = 0.72$ times the CI value:
28 per cent of the chondritic silicon would be missing.
\end{solution}

\begin{solution}{pb:b2:elements-in-the-earth:1}
\textbf{1.} Fe/Mg $= \num{182800}/\num{95870} = 1.907$; Ni/Mg $= 0.1110$; Si/Mg $=
1.111$; Al/Mg $= 0.0887$.

\textbf{2.} $(\num{182800}/55.85)/(\num{106500}/28.09) = 3273/3791 = 0.863$ iron atom per
silicon atom, within 7 per cent of the photosphere's 0.81: sensible, since \omterm{def:b2:elements-in-the-earth:ci}{CI chondrites}
match the Sun for the rock-forming elements.

\textbf{3.} The loss of a component that holds neither element of a ratio leaves the
ratio unchanged (\cref{exo:b2:elements-in-the-earth:12}): only the absolute concentrations change.

\textbf{4.} Magnesium is lithophile and does not enter the core; it is a major element,
well determined in both the \omterm{def:b2:planet-formation-meteorites:chondrule}{chondrites} and the mantle. Iron is the element being shared
between core and mantle: it cannot be its own reference.

\textbf{5.} Elements that condense together were not separated by volatility in the
nebula, so the bulk Earth should keep their ratios. Al/Mg departs: aluminium is
refractory (\qty{1653}{K}), and the silicate Earth's Al/Mg is 1.16 times the CI value.

\textbf{6.} Fe $= 8.05\times55.85/71.85 = 6.26$ per cent; Ni $= 0.25\times58.69/74.69 =
0.196$ per cent.

\textbf{7.} Fe/Mg $= 6.26/22.8 = 0.2746$; Ni/Mg $= 0.196/22.8 = 0.00860$.

\textbf{8.} Iron: $1 - 0.2746/1.907 = 0.856$, 85.6 per cent missing; nickel:
$1 - 0.00860/0.1110 = 0.923$, 92.3 per cent.

\textbf{9.} Iridium is more strongly siderophile than nickel, and nickel than iron:
the more an element prefers metal, the more of it the metal took to the core.

\textbf{10.} In the silicates, as \ce{Fe^2+} in olivine and pyroxene (8 per cent FeO):
iron oxidised to FeO bonds with oxygen like a \omterm{def:b2:elements-in-the-earth:goldschmidt}{lithophile element} and stayed in the
mantle; only the metallic iron sank.

\textbf{11.} $(1-x)C_{\mathrm{Fe}} + xc_{\mathrm{Fe}} = r_{\mathrm{Fe}}(1-x)C_{\mathrm{Mg}}$;
dividing by $(1-x)c_{\mathrm{Fe}}$ gives the proposition.

\textbf{12.} $r_{\mathrm{Fe}}C_{\mathrm{Mg}} - C_{\mathrm{Fe}} = 1.907\times22.8 - 6.26 =
37.22$; $y = 37.22/85.5 = 0.4353$, $x = y/(1+y) = 0.303$.

\textbf{13.} $0.1110\times22.8 - 0.196 = 2.335$; $y = 2.335/5.2 = 0.4490$, $x = 0.310$.

\textbf{14.} $y = (85.5\times37.22 + 5.2\times2.335)/(85.5^2 + 5.2^2) =
3194.4/7337.3 = 0.4354$, $x = 0.303$.

\textbf{15.} Each equation weighs as the square of its coefficient: $85.5^2/5.2^2 = 270$
times more for iron, the far more abundant element.

\textbf{16.} $37.22/2.335 = 15.9$, against $85.5/5.2 = 16.4$ in the model: 3 per cent
apart.

\textbf{17.} $0.303\times\num{5.972e24} = \qty{1.81e24}{kg}$.

\textbf{18.} $(32.5-30.3)/32.5 = 7$ per cent lower.

\textbf{19.} $y = 0.325/0.675 = 0.4815$, Fe in the core $= 37.22/0.4815 = 77.3$ per cent,
and Ni $= 2.335/0.4815 = 4.85$ per cent: 18 per cent light elements, about twice the 5 to
10 per cent allowed by the density. The Earth's Fe/Mg must be somewhat higher than the CI
value, or its silicate magnesium higher than the model's.

\textbf{20.} Silicon, sulfur, carbon, oxygen and hydrogen. Silicon is short in the
silicate Earth: its Mg-normalised ratio, 0.83, is below the refractory level.

\textbf{21.} $1.111\times22.8 - 21.0 = 4.331$; $4.331/0.4354 = 9.9$ per cent Si in the
core.

\textbf{22.} 9.9 against 6 per cent: the core may hold more silicon than the model, or
part of the silicon was never accreted. Silicon condenses at \qty{1310}{K}, a little
below magnesium and iron (1336 and \qty{1334}{K}): some may have stayed in the gas.

\textbf{23.} Core: $0.325\times0.5 = \qty{0.1625}{ppm}$ of the Earth's mass; silicate
Earth: $0.675\times0.001 = \qty{0.000675}{ppm}$. The core holds
$0.1625/0.1632 = 99.6$ per cent of the gold.

\textbf{24.} Chemistry alone puts
\textbf{\qty[text-series-to-math=true]{30.3}{\percent}} of the Earth's mass in its core,
within 7 per cent of the seismological 32.5 per cent.
\end{solution}
